% Chapter 12: The Claude Maturity Model
\chapter{The Claude Maturity Model}
\label{ch:maturity}

Adopting Claude Code is not a binary event.
Teams progress through distinct levels of sophistication,
each building on the previous. This maturity model provides
concrete milestones so you can assess where you are
and what to invest in next.

\section{Five Levels of Claude Adoption}
\label{sec:maturity-levels}

\begin{fullwidth}
\begin{tabular}{@{}clp{4cm}p{5cm}@{}}
\toprule
\textbf{Level} & \textbf{Name} & \textbf{Key Capabilities} & \textbf{Milestone} \\
\midrule
\maturitylevel{1} & Conversational & Default Claude, chat-style usage, no configuration &
  Create first CLAUDE.md (\code{/init}) \\[6pt]
\maturitylevel{2} & Configured & CLAUDE.md, custom commands, permissions, basic hooks &
  5+ commands in a workflow chain \\[6pt]
\maturitylevel{3} & Systematic & Hub-and-spoke architecture, conditional rules, sessions, memory, plan documents &
  Hub repo with 3+ spoke projects \\[6pt]
\maturitylevel{4} & Autonomous & Skills, MCP servers, subagents, parallel workflows, quality dashboards &
  Autonomous quality pipeline with traffic-light dashboard \\[6pt]
\maturitylevel{5} & Enterprise & Agent SDK, agent teams, CI/CD integration, team-wide standards, SSO &
  Team-wide deployment with usage metrics \\
\bottomrule
\end{tabular}
\end{fullwidth}

% -------------------------------------------------------------------
\section{Level Descriptions}

\subsection{L1: Conversational (Week 1)}

Most developers start here: using Claude Code as a chatbot with file access.
No CLAUDE.md, no configuration, no persistence between sessions.

\textbf{Typical activities:} Ask questions, generate code snippets,
explain existing code.

\textbf{Graduation criteria:} Realize that Claude keeps asking the same
questions about your project (build commands, conventions, file locations)
and create a CLAUDE.md to answer them once.

\subsection{L2: Configured (Weeks 2--4)}

The project has a CLAUDE.md with build commands and code conventions.
Custom commands encode common workflows. Permissions prevent unsafe operations.
Basic hooks auto-format code.

\textbf{Typical activities:} Guided development with project context,
custom command workflows, permission-protected operations.

\textbf{Graduation criteria:} Multiple projects share common patterns,
and maintaining them independently becomes a burden.

\subsection{L3: Systematic (Months 2--3)}

\margintip{L3 is where most solo developers should stabilize. L4--L5 add value primarily for teams and complex projects.}

A central hub repo holds shared patterns referenced by spoke projects.
Conditional rules load context based on which files are being edited.
Sessions persist across days with memory retaining cross-session decisions.
Large tasks use plan documents with lifecycle management.

\textbf{Typical activities:} Multi-project workflows with shared standards,
plan-driven development, cross-session memory.

\textbf{Graduation criteria:} Repeated multi-step workflows would benefit
from automation (skills), or external tool access would save significant time.

\subsection{L4: Autonomous (Months 3--6)}

Skills automate multi-step workflows. MCP servers connect Claude to external
systems. Subagents handle isolated tasks in parallel. Quality dashboards
run autonomously and surface issues proactively.

\textbf{Typical activities:} Autonomous quality pipelines, parallel code review,
cross-system queries, metrics-driven improvement sprints.

\textbf{Graduation criteria:} Team adoption requires standardized configuration,
CI/CD integration, or compliance with organizational policies.

\subsection{L5: Enterprise (Month 6+)}

The Claude Agent SDK powers production automation.
Agent teams collaborate on complex tasks.
CI/CD pipelines use Claude for automated review and testing.
Team-wide standards are enforced via managed settings.
SSO and audit logs satisfy compliance requirements.

\textbf{Typical activities:} Production automation, team-wide deployments,
compliance-ready workflows, usage analytics.

% -------------------------------------------------------------------
\section{Progression Strategy}
\label{sec:progression}

\practitioner{Promote when friction justifies it, not on a schedule.}

\begin{enumerate}
  \item \textbf{Stay at L1} until you find yourself repeating the same
    context to Claude across sessions.
  \item \textbf{Move to L2} when a CLAUDE.md and 3--5 commands would
    eliminate daily repetition.
  \item \textbf{Move to L3} when you have 3+ projects sharing patterns
    or when plan documents would prevent wasted effort on large tasks.
  \item \textbf{Move to L4} when repeated multi-step workflows consume
    significant time, or when external tool access would unlock new capabilities.
  \item \textbf{Move to L5} when team adoption, compliance, or production
    automation requires enterprise infrastructure.
\end{enumerate}

\begin{warnbox}[Premature Promotion]
Jumping from L1 to L4 is a common mistake.
Each level builds on the habits and infrastructure of the previous one.
Skipping levels means the foundation is missing:
skills without CLAUDE.md produce inconsistent results,
MCP servers without permissions create security risks,
agent teams without plan documents produce uncoordinated work.
\end{warnbox}

% -------------------------------------------------------------------
\section{Visual: Maturity Pyramid}

\begin{figure}[H]
\centering
\begin{tikzpicture}[
  level/.style={draw=PrimaryBlue, fill=PrimaryBlue!#1, trapezium, trapezium angle=80,
    minimum height=0.8cm, font=\small, align=center, inner sep=2pt},
  >=Stealth
]
  % Pyramid levels (bottom to top, wider to narrower)
  \node[level=5, minimum width=11cm] (l1) at (0,0) {\textbf{L1 Conversational} --- Default usage, no config};
  \node[level=8, minimum width=9cm] (l2) at (0,1.1) {\textbf{L2 Configured} --- CLAUDE.md, commands, permissions};
  \node[level=12, minimum width=7cm] (l3) at (0,2.2) {\textbf{L3 Systematic} --- Hub-and-spoke, rules, memory};
  \node[level=16, minimum width=5cm] (l4) at (0,3.3) {\textbf{L4 Autonomous} --- Skills, MCP, subagents};
  \node[level=22, minimum width=3.2cm, text=white] (l5) at (0,4.4) {\textbf{L5 Enterprise}};

  % Timeline annotations
  \node[font=\scriptsize\itshape, text=PrimaryBlue!70, anchor=west] at (5.8,0) {Week 1};
  \node[font=\scriptsize\itshape, text=PrimaryBlue!70, anchor=west] at (4.8,1.1) {Weeks 2--4};
  \node[font=\scriptsize\itshape, text=PrimaryBlue!70, anchor=west] at (3.8,2.2) {Months 2--3};
  \node[font=\scriptsize\itshape, text=PrimaryBlue!70, anchor=west] at (2.8,3.3) {Months 3--6};
  \node[font=\scriptsize\itshape, text=PrimaryBlue!70, anchor=west] at (1.9,4.4) {Month 6+};

  % Stabilization line
  \draw[dashed, thick, TipGreen!70!black] (-5.5,2.7) -- (5.5,2.7);
  \node[font=\scriptsize\bfseries, text=TipGreen!70!black, anchor=east] at (-5.5,2.7) {Most solo devs};
\end{tikzpicture}
\caption{The maturity pyramid. Most solo developers stabilize at L3.}
\label{fig:maturity-pyramid}
\end{figure}

\begin{keyinsight}
The maturity model is not aspirational---it is diagnostic.
Use it to identify your current level, understand what is holding
you back, and invest in the specific capability that would
produce the most value for your next level.
\end{keyinsight}
